% =============================================================================
% Related Work v2 --- expanded related work for NeurIPS 2026 submission
% Replaces the Related Work section in paper/main.tex (sec:related).
% Grouped into five themes:
%   (1) RL algorithms for LLMs
%   (2) Reasoning RL (and process reward models)
%   (3) Tool-use RL
%   (4) Scaling laws for RL post-training
%   (5) Frameworks, benchmarks, and reproducibility
% Every paragraph references entries defined in paper/references.bib.
% =============================================================================

\section{Related Work}
\label{sec:related}

\ours{} sits at the intersection of five rapidly evolving lines of work:
policy-gradient \emph{algorithms} for language models, RL for \emph{reasoning}
(including process reward models), RL for \emph{tool use} and agentic behaviour,
empirical \emph{scaling laws} for RL post-training, and the \emph{frameworks}
and \emph{benchmarks} that instantiate these algorithms. The closest diagnostic and methodological precedents are distinguished below
from the broader literature context.

% -----------------------------------------------------------------------------
\subsection{RL Algorithms for Large Language Models}
\label{sec:related:algos}

Proximal Policy Optimization (PPO)~\cite{schulman2017proximal} is the workhorse
algorithm behind InstructGPT-style RLHF~\cite{ouyang2022training, christiano2017deeprlhf, stiennon2020summarize, bai2022hhrlhf, bai2022claudecai}
and has dominated production deployments for five years. Whether
group-relative methods beat it is stack-dependent: careful critic-based
variants remain competitive~\cite{kazemnejad2025vineppo, liu2025asyppo}, and
controlled multi-algorithm studies on shared infrastructure find
scale-dependent ranking inversions~\cite{li2026oxrl}; we therefore treat any
single-stack GRPO-vs-PPO gap as a replication-level observation. Direct Preference
Optimization~\cite{rafailov2024direct} reframed preference learning as
contrastive classification and spawned a large family of RL-free alternatives,
including IPO~\cite{azar2024ipo}, KTO~\cite{ethayarajh2024kto},
SimPO~\cite{meng2024simpo}, ORPO~\cite{hong2024orpo}, and
SLiC~\cite{zhao2023slic}, which are possible no-rollout baselines for future work, not evaluated
comparators in this audit.

The current wave is defined by \emph{group-relative} policy gradient methods.
GRPO was introduced with DeepSeekMath~\cite{shao2024deepseekmath} and
popularised by DeepSeek-R1~\cite{guo2025deepseekr1, guo2025deepseekr1nature}, whose R1-Zero study showed that
pure outcome-reward RL can elicit chain-of-thought behaviour from a base
model without supervised cold-start data. A second wave of refinements
dissects the biases of vanilla GRPO. Dr.~GRPO~\cite{tong2025drgrpo} removes a
response-length bias and a question-difficulty amplification bias by using a
constant loss normaliser and dropping the standard-deviation denominator in the
advantage. MC-GRPO (median-centered advantage)~\cite{mcgrpo2025} replaces the
group-mean baseline in the advantage with the group median to stabilise
group-relative normalisation under small rollout counts. G$^{2}$RPO~\cite{hu2026openvlthinker} forces per-task
advantage distributions to match $\mathcal{N}(0,1)$ to fix reward-topology variance in
multimodal settings. Wu et al.~\cite{wu2025grpodpo} relate GRPO and DPO through contrastive-gradient
representations, with differences in token/group weighting and regularization.
Group size can change the expected update, not only Monte Carlo variance.
Their reported efficiency experiments do not establish a saving in our corpus.

Recent 2026 work pre-empts this diagnostic. The Advantage Collapse Rate
(ACR) of AVSPO~\cite{he2026advantage} measures the same homogeneous-group
advantage collapse that we call ZVF (ACR thresholds reward SD while ZVF uses exact equality; matching thresholds
and prompt weights makes the basic measured object coincide) and reports that it predicts stagnation and final
performance; Gradient Starvation~\cite{nie2026gradstarvation} proves that the true
degeneracy rate exceeds the i.i.d.\ Bernoulli prediction by Jensen's
inequality; DARA~\cite{zheng2026dara} reweights rewards by their active-group
density in multi-reward RL; VCRL~\cite{jiang2025vcrl} and online difficulty
filtering~\cite{bae2025online} select prompts by group reward variance or
intermediate pass rate; query recycling treats zero-variance queries
as candidates for later resampling~\cite{coelho2026recycling}; and ConSPO
analyzes a clipped, contrastive GRPO objective and proposes a distinct
sequence-level alternative~\cite{zhang2026conspo}. Goldilocks RL instead
adapts task difficulty under a GRPO training setup~\cite{mahrooghi2026goldilocks}.
These works are prior or concurrent interventions, not evidence for our reviewed
runner. Accordingly, this paper claims neither priority nor an intervention
effect: it confines itself to a submitted-runner accounting identity and an
evidence audit.

\begin{table}[htbp]
\centering\small
\caption{Closest diagnostic comparison; external validation does not transfer
without a matched evaluation.}
\label{tab:diagnostic-prior}
\begin{tabular}{p{.17\linewidth}p{.34\linewidth}p{.38\linewidth}}
\toprule
Quantity & Definition and assumptions & Evidence boundary \\
\midrule
ZVF (this audit) & Exact equality within each sampled group; equal prompt weights & Two descriptive all-wrong cells; no validated operator policy \\
ACR~\citep{he2026advantage} & Fraction with group reward SD below a specified tolerance & External diagnosis/mitigation experiments; no priority claim here \\
$h_G(p)$ & Homogeneity probability for conditionally iid binary draws at fixed prompt success $p$ & Mathematical reference, not observed ZVF or an aggregate-reward predictor \\
\bottomrule
\end{tabular}
\end{table}

A parallel thread focuses on \emph{production-scale} refinements. GSPO, from
the Qwen team~\cite{qwen2025gspo}, redefines the importance ratio at the
sequence level (rather than token level) and was credited with stabilising the
Qwen3 series, especially its Mixture-of-Experts variants. DAPO, from
ByteDance~\cite{yu2025dapo}, contributes asymmetric Clip-Higher, dynamic
sampling, token-level loss normalisation, and overlong filtering, reaching
$50$ points on AIME~2024 with a 32B model. Kimi-K1.5~\cite{team2025kimi15, team2025kimik2} and DeepSeek-V3~\cite{deepseek2024v3} both report their own
GRPO variants and online mirror-descent-style updates. Other policy-gradient work includes VAPO~\cite{bytedance2025vapo}.
Ahmadian et al.~\cite{ahmadian2024backtobasics} revisit REINFORCE-style
optimization, including a leave-one-out baseline, against PPO, DPO, and RAFT;
that contribution is not ``Dr.~SAC'' or evidence of matching GRPO.

A third thread attacks the \emph{zero-variance failure mode} (ZVF) directly.
GRESO~\cite{zhang2026greso} pre-screens prompts
before sampling; CPPO~\cite{lin2025cppo} prunes completions post-rollout,
giving up to $7.98\times$ speedup on GSM8K. NGRPO~\cite{nan2025ngrpo} adds
Advantage Calibration and asymmetric clipping so that all-incorrect groups
still carry gradient. Scaf-GRPO~\cite{zhang2025scaffgrpo} injects tiered
scaffolding hints when learning stagnates, boosting AIME~2024 by $44.3\%$
relative. MO-GRPO~\cite{ichihara2025mogrpo} addresses vanilla GRPO's
collapse to a single reward component under multiple objectives via
variance-based reward re-weighting. Finally,
SSPO~\cite{yang2025sspo} computes importance ratios at the subsentence level
to balance token-level GRPO against sequence-level GSPO, reporting an average
of $46.72$ versus $43.01$ (GRPO) and $44.42$ (GSPO) over five math datasets
on Qwen2.5-1.5B-Math. Taken together, these results explain at
a theoretical level why implementation details at the granularity of the
importance ratio have such large empirical effects --- a motivation for auditing objective scope and cohort identity; this audit
does not identify implementation defaults as the cause of a measured gap.

% -----------------------------------------------------------------------------
\subsection{Reasoning RL and Process Reward Models}
\label{sec:related:reasoning}

The shift from outcome rewards to \emph{process} supervision has reshaped
reasoning RL. Let's-Verify-Step-by-Step~\cite{lightman2023prm} showed
that step-level PRMs outperform outcome reward models on MATH, following the
GSM8K process- versus outcome-feedback comparison of
Uesato et al.~\cite{uesato2022prm800k}. Math-Shepherd~\cite{wang2024mathshepherd} derives process rewards
from auto-generated process annotations, avoiding costly human labels.
PRM800K~\cite{lightman2023prm} and PRMBench~\cite{song2024prmbench} provide
evaluation data for step-level scoring. OmegaPRM~\cite{luo2024omegaprm} uses
Monte-Carlo tree search to supervise process rewards without human-in-the-loop,
while ReasonEval~\cite{xia2024reasoneval} scores the validity and redundancy
of reasoning steps beyond final-answer accuracy. Implicit PRM~\cite{yuan2024implicitprm} learns PRMs implicitly from
outcome labels, and VersaPRM~\cite{yuan2025versaprm} extends process rewards
to multi-domain tasks. Process-reward approaches must be distinguished from outcome/format reward
training. In particular, DeepSeek-R1 discusses process reward models among its
unsuccessful attempts~\cite{guo2025deepseekr1}. We do not attribute a specific
reward algorithm to proprietary reasoning systems without disclosed methods.

Complementary results demonstrate that
self-consistency~\cite{wang2023selfconsistency}, tree-of-thought
search~\cite{yao2024tot}, and critic-style post-hoc
verifiers~\cite{mcaleese2024criticgpt} all benefit from or substitute for
process reward supervision. VerifyBench~\cite{zhang2025verifybench} and
A~Sober Look~\cite{hochlehnert2025sober} warn, respectively, that verifier
errors and evaluation artefacts (decoding, seeds, prompt format, hardware and
software configuration) can inflate perceived reasoning gains, motivating our
insistence on binary verifiable rewards and reported evaluation stacks.

Reasoning-specific algorithmic contributions also continue to accumulate.
ReFT~\cite{luong2024reft} studies supervised fine-tuning versus RL for
reasoning. Let's Reward Step~\cite{zhang2025rewardstep} uses step-level PRM
feedback to guide greedy search at inference time, and
STaR~\cite{zelikman2022star} bootstraps reasoning chains from its own
correct rationales. rStar-Math~\cite{guan2025rstarmath} combines MCTS-based PRMs
with self-evolved reasoning. ReST-EM~\cite{singh2024restem} alternates
expectation and maximisation steps over self-generated rationales.
Self-Taught Evaluator~\cite{wang2024stevaluator} adapts an LLM to score its
own reasoning chains. GRPO-Lead~\cite{lu2025grpolead}, LCPO~\cite{lu2025lcpo},
and LIMR~\cite{li2025limr} each tune reasoning depth with length-aware
rewards. Search-R1~\cite{jin2025searchr1} trains reasoning agents to
interleave search queries with chain-of-thought. TTRL~\cite{pu2025ttrl}
performs test-time RL on unlabeled prompts using majority-vote pseudo-labels
as rewards. Dr.~GRPO~\cite{tong2025drgrpo} and GSPO~\cite{qwen2025gspo} are external
formulations; this audit does not evaluate them as reference implementations.

% -----------------------------------------------------------------------------
\subsection{Tool-Use and Agentic RL}
\label{sec:related:tooluse}

Tool-use RL trains LLMs to interleave natural-language reasoning with external
API calls. Toolformer~\cite{schick2023toolformer} introduced self-supervised
tool-call insertion for a small set of APIs; Gorilla~\cite{patil2023gorilla}
retrieved from a large API catalogue; and ToolLLM~\cite{qin2024toolllm} scaled
training to thousands of real-world REST APIs. ReAct~\cite{yao2023react}
established the interleaved reason--act--observe pattern that most subsequent
agents inherit. ToolACE~\cite{liu2024toolace} generates function-calling data and evaluates
supervised fine-tuning, rather than establishing a tool-use RL result; ToolRL~\cite{qian2025toolrl} applies GRPO
with verifiable tool-call rewards; ToRL~\cite{li2025torl} optimises the joint
reasoning--action trajectory. Further advances include
APIGen~\cite{lin2024apigen} (high-quality synthetic API data),
Nexus-Raven~\cite{nexusflow2024nexusraven}, and
Octopus~\cite{chen2024octopus}. Agentic RL extends tool use to long-horizon
tasks: WebGPT~\cite{nakano2021webgpt} and
WebShop~\cite{yao2022webshop} were the first large-scale web agents,
Reflexion~\cite{shinn2023reflexion} and Voyager~\cite{wang2023voyager} add
episodic self-reflection, and AutoGPT-style agents have been formalised by
AgentBench~\cite{liu2024agentbench}. 2025--2026 produced a wave of RL-trained
agents: SWE-RL~\cite{wei2025swerl} trains coding agents with execution
rewards; SWE-Gym~\cite{pan2025swegym} provides an RL training environment for
real GitHub issues; ARTIST~\cite{singh2025artist} trains tool-integrated
reasoning agents with GRPO; RAGEN~\cite{wang2025ragen} studies multi-turn
RL with tool access; OpenAI's Deep Research~\cite{openai2025deepresearch}
and Perplexity Deep Research~\cite{perplexity2025deepresearch} motivate
web-agent evaluation; no specific GRPO implementation is attributed here. Our benchmark currently focuses
on single-turn verifiable mathematical reasoning; we adopt Tinker's
\texttt{message\_with\_tool} interface as the neutral substrate for future
tool-use extensions and explicitly cite these works as concurrent settings in
which \ours{}'s ZVF measurements will need to be re-evaluated.

% -----------------------------------------------------------------------------
\subsection{Scaling Laws for RL Post-Training}
\label{sec:related:scaling}

Compute-optimal training of base language models is governed by
well-known scaling laws~\cite{kaplan2020scalinglaws, hoffmann2022chinchilla, muennighoff2023datascaling}, but RL post-training is much less well
characterised. Hilton et al.~\cite{hilton2023rlscaling} first derived
compute-efficiency frontiers for RL in games and robotics. For RLHF,
Gao et al.~\cite{gao2023reward} characterise the over-optimisation curve of
reward models, and Rafailov et al.~\cite{rafailov2024rlhfscaling} provide
scaling laws for DPO and PPO that show a strong dependence on reward-model
quality. Nimmaturi et al.~\cite{nimmaturi2025predictive} derive an
empirical scaling law specifically for GRPO, identifying three training
phases (slow start, rapid improvement, plateau) and showing that most
benefit is exhausted by roughly $80\%$ of one epoch. Mukherjee et
al.~\cite{liu2025rlsubnetworks} find that RL fine-tuning updates only
$5$--$30\%$ of model parameters across seven algorithms and ten model
families, suggesting RL activates a sparse ``RL subnetwork.''
MiniCPM~\cite{hu2024minicpm} and the SLM survey of
Nguyen et al.~\cite{subramanian2025slm} characterise the small-model
regime. Schaeffer et
al.~\cite{schaeffer2024emergent} caution that apparent emergent abilities may
be metric artefacts, a concern that persists into RL. Our short heterogeneous traces are descriptive context, not scaling-law
ground truth or independent scale replications.

% -----------------------------------------------------------------------------
\subsection{Frameworks, Benchmarks, and Reproducibility}
\label{sec:related:frameworks}

The proliferation of algorithms has been matched by a proliferation of
training \emph{frameworks}: TRL~\cite{vonwerra2022trl},
OpenRLHF~\cite{hu2024openrlhf}, veRL~\cite{sheng2024verl},
NeMo-RL~\cite{nvidia2024nemorl}, DeepSpeed-Chat~\cite{yao2023deepspeedchat},
trlX~\cite{havrilla2023trlx}, Axolotl~\cite{axolotl2024}, and
\tinker{}~\cite{thinkingmachines2024tinker}. Each framework makes different
default choices about advantage normalisation, KL estimators, tokenisation
loss masking, rollout batching, and importance-sampling granularity; those
differences are possible confounders, not identified causes in this audit.
On the inference side, vLLM~\cite{kwon2023vllm} and
SGLang~\cite{zheng2024sglang} provide the rollout backends on which all
modern frameworks depend; FlashAttention~\cite{dao2022flashattention} and
LoRA~\cite{hu2022lora} are load-bearing kernels.

In the RL \emph{benchmarking} tradition, Henderson et
al.~\cite{henderson2018deep} documented the fragility of deep-RL results
to implementation choices and seeds;
\texttt{rliable}~\cite{agarwal2021deep} introduced interquartile-mean
reporting and performance profiles; BenchMARL~\cite{bettini2023benchmarl}
unified multi-agent RL benchmarks; Patterson et
al.~\cite{patterson2024empirical} surveyed empirical design practices;
Jordan et al.~\cite{jordan2024benchmarking} argued that most published RL
comparisons lack statistical power; and Madaan et
al.~\cite{zhang2024benchvariance} quantified LLM evaluation variance along
seed, monotonicity, and scaling axes. Pineau et
al.~\cite{pineau2020improving} formalised NeurIPS reproducibility criteria
and ACM's Artifact Review and Badging~\cite{acm2020badges} established
community standards for available, evaluated, and reproduced artifacts.
Reproducibility studies in adjacent fields include
\cite{cavenaghi2023systematic, lynnerup2019survey, yamerenko2025generating, paparella2023reproducibility, huang2022state, lopezibanez2021reproducibility, leeser2026artifact, sedghpour2024artifact, colas2019hitchhiker}, alongside
Hoffman et al.'s Acme~\cite{hoffman2022acme}. Classical evaluations rest on
GSM8K~\cite{cobbe2021training}, MATH~\cite{hendrycks2021math},
AIME~\cite{maa2024aime}, MMLU~\cite{hendrycks2021mmlu},
HumanEval~\cite{chen2021codex}, MBPP~\cite{austin2021mbpp}, and
BBH~\cite{suzgun2023bbh}. Implementation Matters~\cite{engstrom2020implementation}
directly examines how implementation choices affect apparent PPO/TRPO gains.
Against this prior work, our narrower contribution is an objective-scoped
identity, an evidence audit, and a reporting contract. We claim neither a first
cross-library protocol nor novelty for the ZVF definition.

% =============================================================================
